\documentclass[a4,11pt, oneside]{book}
\usepackage[utf8]{inputenc}
\usepackage{lmodern}
%\frenchbsetup{StandardLists=true}
\usepackage[usenames]{color}
\usepackage{latexsym}
\usepackage{amssymb}
\usepackage{textcomp}
\usepackage{multirow}
\usepackage{amsmath, amsfonts, amssymb}
\usepackage{enumitem}
\usepackage{mathrsfs}% \mathscr{C}
\usepackage[T1]{fontenc}
\usepackage[french]{babel}
%\usepackage[french, english]{babel}
\usepackage{graphicx}
%\usepackage{slashbox}
\usepackage{a4wide}
\usepackage{lscape}
\addtocounter{chapter}{0}
\usepackage{rotating}
\usepackage{fancyhdr}
\usepackage[a4paper]{geometry}
% \usepackage{titlesec}
\usepackage[colorlinks=false,urlcolor=black]{hyperref} % Actionner les liens hypertextes
%%%%%%%%%% Mise en forme de la table des matieres %%%%%%%
\newcommand{\chaptertoccolor}{black} %couleur des chapitres
\newcommand{\sectiontoccolor}{blue} % couleur des sections
\newcommand{\subsectiontoccolor}{black} % couleur des sous-sections {green}
\newcommand{\subsubsectiontoccolor}{red}
\setcounter{tocdepth}{2} % definit la profonduer de la table des matieres
\makeatletter
% On redéfinit la façon dont doit être écrit le titre des chapitres
\renewcommand*\l@chapter[2]{
	\ifnum \c@tocdepth >\m@ne
	\addpenalty{-\@highpenalty}
	\vskip 1.0em \@plus\p@
	\setlength\@tempdima{1.5em}
	\begingroup
	\parindent \z@ \rightskip \@pnumwidth
	\parfillskip -\@pnumwidth
	\leavevmode \bfseries
	\advance\leftskip\@tempdima
	\hskip -\leftskip
	\def\@linkcolor{\chaptertoccolor}
	\color{\chaptertoccolor}{\mathversion{bold} #1}\nobreak\
	\leaders\hbox{$\m@th
		\mkern \@dotsep mu\hbox{.}\mkern \@dotsep
		mu$}\hfil\nobreak\hb@xt@\@pnumwidth{\hss
		\def\@linkcolor{\chaptertoccolor}
		\color{\chaptertoccolor}#2}\par
	\penalty\@highpenalty
	\endgroup
	\fi}
% Définition des couleurs des différents niveaux dans la table des matières
\renewcommand*\l@section{\color{\sectiontoccolor}
	\def\@linkcolor{\sectiontoccolor}
	\@dottedtocline{1}{1.5em}{2.3em}}
\renewcommand*\l@subsection{\color{\subsectiontoccolor}
	\def\@linkcolor{\subsectiontoccolor}
	\@dottedtocline{2}{2.5em}{3.3em}}
\renewcommand*\l@subsubsection{\color{\subsubsectiontoccolor}
	\def\@linkcolor{\subsubsectiontoccolor}
	\@dottedtocline{3}{3.5em}{4.3em}}

\makeatother
%%%%% fin de la mise en forme de la table des matieres %%%%%%%%%

%%%%%%%%%%%%%%%%%%%%%%%%%%
%\usepackage{hyperref}



\usepackage{pxfonts}
%%%diagrams packages xy$
\usepackage{xy}
% \input{./texinputs/xy.sty}
% \input xy
\xyoption{all}
%%%%%%%%%%%%%%%
\usepackage{color}
%%%%%%%%%%%%%%%%%%%%%%%% This show the scale of a page%%%%%%%%%%%%%%%%%%%%%%
%%%%%%%%%%%%%%%%%%%%%%%%%%%%%%%%%%%%%%%%%%%%%%%%%%%%%%%%%%%%%%%%%%%%%%%%%%%
%\geometry{hscale=0.70,vscale=0.80}
%%%%%%%%% This numbers equations by sections %%%%%%%%%%%%%%%%%%%%%%%%%%%%%
%%%%%%%%%%%%%%%%%%%%%%%%%%%%%%%%%%%%%%%%%%%%%%%%%%%%%%%%%%%%%%%%%%%%%%%%%%
% \oddsidemargin -.3cm \evensidemargin 0 cm
%\setlength{\textheight}{22cm}
\textheight 25cm
\textwidth 16cm
%\hoffset -1.7 cm
\voffset -1.5 cm


\makeatletter
\@addtoreset{equation}{section}
\renewcommand{\theequation}{\thesection.\arabic{equation}}
\makeatother
%%%%%%%%% Other definitions and commands %%%%%%%%%%%%%%%%%%%%%%%%%%%%%%%%%
%%%%%%%%%%%%%%%%%%%%%%%%%%%%%%%%%%%%%%%%%%%%%%%%%%%%%%%%%%%%%%%%%%%%%%%%%%
%%%%%%%%%%%%%%%%%%%%%%%%%%%%%%%%%%%%%%%%%%%%%%%%%%%%%%%%%%%%%%%%%%%%%%%%%%

%%%%%%%%%%%%%%%%%%%%%%%%%%%%%%%%%%%%%%%%%%%%%%%%%%%%%%%%%%%%%%%%%%%%%%%%%%
\makeatletter
\def\captionof#1#2{{\def\@captype{#1}#2}}
\makeatother
\newcommand{\subs}[1]{{\vspace*{0.5cm}}%
	{\noindent\underline{#1}}{\addcontentsline{toc}{subsection}{#1}}%
	{\vspace*{0.3cm}}}
%%%%%%%%%%%%%%%%%%%%%%%%%%%%%%%%%%%%%%%%%%%%%%%%%%%%%%%%%%%%%%%%%%%%%%%%%%
%%%%%%%%%%%%%%%%%%%%%%%%%%%%%%%%%%%%%%%%%%%%%%%%%%%%%%%%%%%%%%%%%%%%%%%%%%


%\newtheorem{examples}[subsection]{Exemples}
%\newtheorem{counterexamples}[subsection]{Countre-exemples}
%\newtheorem{counterexample}[subsection]{Countre-examples}
%\newtheorem{properties}[subsection]{Propriétés}
\newtheorem{exercise}{Exercise}[section]
%\newtheorem{remarks}[theorem]{Remarks}
%\newtheorem{definitions}[subsection]{Définitions}
%\newtheorem{proof}{Preuve}[section]
%\newtheorem{prooft}{Proof of theorem}[section]
\usepackage{tikz}

\newcommand{\cqfd}{\hfill $\blacksquare$}
%%%%%%%%%%%%%%%%%%%%%%%%%%%%%%%%%%%%%%%%%%%%%%%%%%%%%%%%%%%%%%%%%%%%%%%%%%
%%%%%%%%%%%%%%%%%%%%%%%%%%%%%%%%%%%%%%%%%%%%%%%%%%%%%%%%%%%%%%%%%%%%%%%%%%

%%%%%%%%%%%%%%%%%%%%%%%%%Notations used in Hopf Algebra %%%%%%%%%%%%%%%%%%
%%%%%%%%%%%%%%%%%%%%%%%%%%%%%%%%%%%%%%%%%%%%%%%%%%%%%%%%%%%%%%%%%%%%%%%%%%
\newcommand{\p}{\frac{p^{\mu-1}}{q^{\nu}}}
\newcommand{\q}{\frac{p^{\mu}}{q^{\nu-1}}}
\newcommand{\f}{\Phi}
\newcommand{\lf}{\left(}
\newcommand{\rg}{\right)}
\newcommand{\da}{\dagger}
%%%%%%%%%%%%%%%%%%%%%%%%%%%%%%%%%%%%%%%%%%%%%%%%%%%%%%%%%%%%%%%%%%%%%%%%%%
%%%%%%%%%%%%%%%%%%%%%%%%%%%%%%%%%%%%%%%%%%%%%%%%%%%%%%%%%%%%%%%%%%%%%%%%%%

\renewcommand{\bibname}{References}

%%%%%%%%%%%%%%%%%%%%%%%%%%%%%%%%%% voici les note de bas de page %%%%%%%%%%%%%%%%%%%%%%%%%%%%%%%%%%%%%%%%%%%%%%%
\frenchspacing  \linespread{1.1}%1.5avant
%moi%\usepackage{fancyhdr}
%moi%\pagestyle{fancy}
%moi%\usepackage{graphicx}
%\usepackage{t1enc}
\makeatletter
\newcommand{\thechapterwords}{\thechapter}
%{ \ifcase \thechapter\or 1\or 2\or 3\or 4\or 5\or  6\or 7\or 8\or 9\or 10\or 11\fi}
\def\thickhrulefill{\leavevmode \leaders \hrule height 1ex \hfill \kern \z@}
\def\@makechapterhead#1{%
	%\vspace*{50\p@}%
	% \vspace*{15\p@}%
	{\parindent \z@ \centering \reset@font
		\thickhrulefill\quad
		\scshape \@chapapp{} \thechapterwords
		\quad \thickhrulefill
		\par\nobreak
		\vspace*{10\p@}%
		\interlinepenalty\@M
		\hrule
		\vspace*{10\p@}%
		\Huge \bfseries #1\par\nobreak
		\par %\nouppercase{name}
		% m %       \vspace*{15\p@}%
		\hrule
		% m %   \vskip 60\p@
		%\vskip 100\p@
}}
\def\@makeschapterhead#1{%
	%\vspace*{50\p@}%
	% m %  \vspace*{15\p@}%
	{\parindent \z@ \centering \reset@font
		\thickhrulefill
		\par\nobreak
		\vspace*{10\p@}%
		\interlinepenalty\@M
		\hrule
		\vspace*{10\p@}%
		\Huge \bfseries #1\par\nobreak
		\par %\nouppercase{name}
		\vspace*{10\p@}%
		\hrule
		\vskip 40\p@
		%\vskip 100\p@
}}
\def\@makechapterhead#1{%
	%\vspace*{50\p@}%
	\vspace*{15\p@}%
	{\parindent \z@ \centering \reset@font
		\thickhrulefill\quad
		\scshape \@chapapp{} \thechapterwords
		\quad \thickhrulefill
		\par\nobreak
		\vspace*{10\p@}%
		\interlinepenalty\@M
		\hrule
		\vspace*{10\p@}%
		\Huge \bfseries #1\par\nobreak
		\par %\nouppercase{name}
		\vspace*{10\p@}%
		\hrule
		\vskip 40\p@
		%\vskip 100\p@
}}
\frenchspacing \pagestyle{headings}
%%avant c etait \frenchspacing  \linespread{1.1}%1.5avant
%%d\'efinition des ent\^etes
\usepackage{fancyhdr}
\pagestyle{fancy}


\renewcommand{\sectionmark}[1]{\markboth{#1}{}}
\renewcommand{\chaptermark}[1]{\markboth{#1}{}}
%\renewcommand{\sectionmark}[1]{\markright{\thesection\ #1}}
\fancyhf{} % supprime les en-t\^etes et pieds pr\'ed\'efinis
\fancyhead[L,R]{\bfseries\thepage}% Left Even, Right Odd
\fancyhead[L]{\bfseries\rightmark} % Left Odd
\fancyhead[R]{\bfseries\leftmark} % Right Even
\renewcommand{\headrulewidth}{1pt}% filet en haut de page
\addtolength{\headheight}{1pt} % espace pour le filet
\renewcommand{\footrulewidth}{0.5pt}% filet en bas de page
%\addtolength{\footheight}{0.5pt} % espace pour le filet en bas
\fancypagestyle{plain}{ % pages de tetes de chapitre
	\fancyhead{} % supprime l'entete
	%\fancyfoot{} %supprime le pied de page
	\renewcommand{\headrulewidth}{0pt}
}
\newcommand{\clearemptydoublepage}{%
	\newpage{\pagestyle{plain}\cleardoublepage}}

\rhead{{\thepage}} 
\lhead{\textsl{\leftmark}}

%\fancyfoot[RE]{\tiny \textsl{ Domaine(orientation) }}

%  \rhead{\textbf{\thepage}}
% \lhead{\textsl{\leftmark}}


\lhead{\nouppercase{\rightmark}}
% \rhead{\nouppercase{\leftmark}}

% \newcommand{\clearemptydoublepage}{\newpage{\pagestyle{plain}\cleardoublepage}}
% \newcommand{\chapitre}[1]{\chapter{#1}}

%%%%%%% Mise en forme Page de garde
 %%%%%%%%%%%%%%%%%%%%%%%%%
\makeatletter
\newcommand\@cclxvipt{23}
\newcommand\HUGE{\@setfontsize\HUGE\@cclxvipt{50}}
\makeatother
\usepackage{fancybox}

%
\newsavebox{\auteurbm}
\newenvironment{Bonmot}[1]%
{\small\slshape%
	\savebox{\auteurbm}{\upshape\sffamily#1}%
	\begin{flushright}}
	{\\[4pt]\usebox{\auteurbm}
	\end{flushright}\normalsize\upshape}
%

%%%%%%%
%%%%%%%%%%%%%%%%%%%%%%%%%%%%%%%%%%%%%%%%%%%%%%%%%%%%%%%%%%%%%%%%%%%%%%%

%
\newcommand{\Filetdouble}[1]{%
	\newlength{\ecart}
	\settowidth{\ecart}{#1}
	\begin{center}
		\noindent
		\rule[0ex]{\ecart}{1pt} \par
		#1 \par
		\rule[1ex]{\ecart}{1pt}
\end{center}}

 
\everymath{\displaystyle}

\newcommand{\ie}{c'est-à-dire \  }
\newcommand{\ssi}{si et seulement si\  }
\newcommand{\R}{\ensuremath{\mathbb{R}}}
\newcommand{\N}{\ensuremath{\mathbb{N}}}
\newcommand{\C}{\ensuremath{\mathbb{C}}}
\newcommand{\K}{\ensuremath{\mathbb{K}}}
\newcommand{\ep}{{1\leq i \leq p,1\leq j \leq q}}
\newcommand{\eq}{{1\leq i \leq q,1\leq j \leq p}}
\newcommand{\en}{{1\leq i \leq n,1\leq j \leq m}}
\newcommand{\eij}{{1\leq i,j \leq p}}
\newcommand{\eik}{{1\leq i,k \leq q}}
\newcommand{\Degre}{\ensuremath{^\circ}}
\DeclareMathOperator{\sh}{sh}
\DeclareMathOperator{\e}{e}